(1) 令第 \(n\) 胞内两粒子位移为 \(u_n,v_n\)，最近邻弹簧力给 \(m\ddot u_n=K(v_n+v_{n-1}-2u_n)\)，\(M\ddot v_n=K(u_n+u_{n+1}-2v_n)\)。平移一个晶胞后方程保持不变，因此可以用 Bloch 模态而非两个独立振子的解来描述正常传播；弹簧只连接异种质量粒子。
(2) 代入 \(u_n=Ue^{i(nka-\omega t)}\)、\(v_n=Ve^{i(nka-\omega t)}\)，得到行列式 \((2K-m\omega^2)(2K-M\omega^2)-4K^2\cos^2(ka/2)=0\)。解为 \(\omega_\pm^2=K(1/m+1/M)\pm K\sqrt{(1/m+1/M)^2-4\sin^2(ka/2)/(mM)}\)；下支趋于零，是声学支，上支为光学支。
(3) 在 \(k=0\) 时 \(\omega_-^2=0\)、\(\omega_+^2=2K(1/m+1/M)\)；在 \(k=\pi/a\) 时两支分别等于 \(2K/\max(m,M)\) 与 \(2K/\min(m,M)\)。若 \(m\ne M\)，它们之间有非零禁带，无法找到满足实数 \(k\) 的传播频率；当 \(m=M\) 时边界两支恰好相交，所谓双支仅是同一单原子链色散在缩小布里渊区后的折叠。
(4) 对 \(|ka|\ll1\) 展开 \(\sin(ka/2)\simeq ka/2\)，声学支为 \(\omega_-^2\simeq Ka^2k^2/[2(m+M)]\)，故 \(c_s=a\sqrt{K/[2(m+M)]}\)。当 \(m=M\) 时 \(c_s=(a/2)\sqrt{K/m}\)，与原子实际相邻距离 \(a/2\) 的单原子弹簧链一致；若误把晶胞周期当相邻间距，会多出一个错误因子二。
